\documentclass[10pt,a4paper]{amsart}
\usepackage[margin=19mm]{geometry}
\usepackage{amsmath,amssymb,mathtools}
\usepackage[scr=rsfs,cal=euler]{mathalfa}
\usepackage{xfrac}
\usepackage[unicode,hidelinks,bookmarks=false]{hyperref}
\newcommand{\ie}{i.e.\ }
\newcommand{\cf}{cf.\ }
\newcommand{\resp}{respectively\ }
\newcommand{\Subsection}{\subsection}
\providecommand{\nobreakdash}{\nobreak\hbox{-}}
\setlength{\emergencystretch}{2em}
\multlinegap0pt
\usepackage{bm}
\usepackage{bbm}
\usepackage{dsfont}
\usepackage{xspace}
\usepackage{cancel}
\usepackage{mathtools}
\usepackage{amsthm}
\makeatletter
\@ifundefined{theorem}{\newtheorem{theorem}{Theorem}}{}
\makeatother
\title[Quickest Change Detection with Cost-Constrained Experiment D]{Quickest Change Detection with Cost-Constrained Experiment Design}
\author{Patrick Vincent N. Lubenia, Taposh Banerjee}
\date{Source: arXiv:2509.14186v2 (CC BY 4.0)}
\begin{document}
\maketitle

\noindent\emph{Source and licence.} Quickest Change Detection with Cost-Constrained Experiment Design. Version arXiv:2509.14186v2 (CC BY 4.0), distributed under the Creative Commons Attribution 4.0 International licence, as recorded in the arXiv licence metadata for this exact version and shown on the abstract page https://arxiv.org/abs/2509.14186v2. Full rights notice: licenses/arxiv-2509.14186-CC-BY-4.0.txt.\par\smallskip

\noindent\textit{Editorial scope.} Counted unit: the Theorem whose source label is \texttt{thm:3e} in the article (source0.tex lines 1195--1222, source label \texttt{thm:3e}), delivered together with the article's own complete original proof of it (source0.tex lines 1224--1273, one \texttt{proof} environment of 50 lines). No mathematics is written, repaired or re-derived: statement and proof are the article's own text line for line. Required context retained uncounted: none. Every definition, display and label of those lines is retained unchanged. Omitted from this package and not redistributed: 1--1194, 1223--1223, 1274--1792. Nothing inside a retained excerpt is omitted or altered: every retained body is delivered exactly as the article's lines are written, so no omission receipt of any kind accompanies this package. Citations: the retained excerpts carry no citation command at all, so no bibliography entry of the article is transcribed and no reference is redistributed. Reference adaptation: none was needed and none was made; every reference inside the delivered unit resolves inside it, so no unresolved reference or citation is produced. Printed slips kept literal: no printed word, symbol, hypothesis or number of the counted statement or of its proof was altered. Layout and class: the article's own class is \texttt{IEEEtran} with packages including cite, graphicx, amsmath, array, stfloats, amssymb, amsthm, hyperref, xcolor; the wrapper is the lane's pinned \texttt{amsart} editorial wrapper at 10pt on A4, which restates the theorem-like environments the retained text needs and adds the article's own macro definitions that the retained text needs (none), with extra wrapper packages (bm, bbm, dsfont, xspace, cancel, mathtools). The delivered numbering is the unit's own. Assets: no retained span contains a figure environment, a graphics-inclusion command, a PDF-page inclusion, a table or a TikZ picture; no image file is copied into this package, the package declares no asset dependency and no image file is redistributed with it. Licence: version arXiv:2509.14186v2 of this article is distributed under the Creative Commons Attribution 4.0 International licence, as recorded in the arXiv licence metadata for this exact version and shown on the abstract page https://arxiv.org/abs/2509.14186v2. A full rights notice is delivered at licenses/arxiv-2509.14186-CC-BY-4.0.txt. Characters: every retained excerpt is pure ASCII, so the delivered engine, XeLaTeX with its default Latin Modern text font, reports no missing character.
\begin{theorem}
\label{thm:3e}

For the \textup{3E-CUSUM} Algorithm, let 
$$ 0 < D(f_1^X \mathrel{\Vert} f_0^X) \leq D(f_1^Y \mathrel{\Vert} f_0^Y) \leq D(f_1^Z \mathrel{\Vert} f_0^Z) < \infty. $$

\begin{enumerate}
\item For any fixed $a_Y > 0$, $a_Z > 0$, $N_X > 0$, and $N_Y > 0$, with $A = \log \gamma$,
$$ \textup{ARLFA} (\Psi_{\textup{3E}}) \geq \gamma. $$

\item For fixed values of $A$, $a_Y > 0$, $a_Z > 0$, $N_X > 0$, and $N_Y > 0$,
\begin{align*}
& \textup{POR}_Z (\tau_{\textup{\fontsize{6 pt}{10 pt}\selectfont 3E}})  = \frac{\mathsf{E}_\infty [\sigma^Z]}{\mathsf{E}_\infty [\sigma^Z] + \mathsf{E}_\infty [\tau_{\textup{\fontsize{6 pt}{10 pt}\selectfont 2E}}(a_Z D_{\sigma^Z}, 0, a_Y, N_X, N_Y)]} \\
& \textup{POR}_Y (\tau_{\textup{\fontsize{6 pt}{10 pt}\selectfont 3E}})  = \frac{\mathsf{E}_\infty [\min
\{ \tau_{\textup{\fontsize{6 pt}{10 pt}\selectfont C}}^Y (\vert a_Z D_{\sigma^Z}\vert), N_Y \}]}{\mathsf{E}_\infty [\sigma^Z] + \mathsf{E}_\infty [\tau_{\textup{\fontsize{6 pt}{10 pt}\selectfont 2E}}(a_Z D_{\sigma^Z}, 0, a_Y, N_X, N_Y)]}.
\end{align*}
The POR expressions are not a function of $A$. Thus, for any $A$, we can always choose values for $a_Y$, $a_Z$, $N_X$, and $N_Y$ to meet any given POR constraint of $\beta_Z$ and $\beta_Y$:
$$ \textup{POR}_Z (\Psi_{\textup{3E}} (A, a_Y, a_Z, N_X, N_Y)) \leq \beta_Z $$
and
$$ \textup{POR}_Y (\Psi_{\textup{3E}} (A, a_Y, a_Z, N_X, N_Y)) \leq \beta_Y. $$

\item For fixed values of $a_Y > 0$, $a_Z > 0$, $N_X > 0$, and $N_Y > 0$, and for each $A$,
$$ \textup{WADD} (\tau_{\textup{\fontsize{6 pt}{10 pt}\selectfont 3E}}) \leq \mathsf{E}_1 [\tau_{\textup{\fontsize{6 pt}{10 pt}\selectfont 3E}}] + N_Y + N_Y N_X. $$
Consequently, 
$$ \textup{WADD} (\Psi_{\textup{3E}}) \leq \textup{WADD} (\tau_{\textup{\fontsize{6 pt}{10 pt}\selectfont C}}^Z) + K_3, $$
where $\tau_{\textup{\fontsize{6 pt}{10 pt}\selectfont C}}^Z$ is the \textup{CUSUM} algorithm for the experiment $Z$ and $K_3$ is a constant that is not a function of threshold $A$. 
\end{enumerate}
\end{theorem}

\begin{proof}
\begin{enumerate}
\item The false alarm analysis for the 3E-CUSUM algorithm is similar to that of the 2E-CUSUM algorithm. Specifically, for any value of threshold $A$, we have
$$ \mathsf{E}_\infty [\tau_{\textup{\fontsize{6 pt}{10 pt}\selectfont 3E}}] \geq \mathsf{E}_\infty [\tau_{\textup{\fontsize{6 pt}{10 pt}\selectfont C}}^Z (A)] \geq e^A. $$

\item The expression for $\textup{POR}_Z$ follows again from the Renewal Reward Theorem because under the pre-change distribution, the expected number of observations of experiment $Z$ is $\mathsf{E}_\infty [\sigma^Z]$. Also, the average time spent below zero is given by the time for the truncated 2E-CUSUM algorithm to stop and bring the statistic above $0$: $\mathsf{E}_\infty [\tau_{\textup{\fontsize{6 pt}{10 pt}\selectfont 2E}}(a_Z D_{\sigma^Z}, 0, a_Y, N_X, N_Y)]$. This gives
\begin{align*}
\textup{POR}_Z (\tau_{\textup{\fontsize{6 pt}{10 pt}\selectfont 3E}})  = \frac{\mathsf{E}_\infty [\sigma^Z]}{\mathsf{E}_\infty [\sigma^Z] + \mathsf{E}_\infty [\tau_{\textup{\fontsize{6 pt}{10 pt}\selectfont 2E}}(a_Z D_{\sigma^Z}, 0, a_Y, N_X, N_Y)]}.
\end{align*}

When $N_Y = 0$, i.e., experiment $Y$ is never performed (consequently, experiment $X$ is never performed as well), $\textup{POR}_Z = 1$ and experiment $Z$ is utilized all the time. To decrease the fraction of time experiment $Z$ is chosen, we increase $N_Y$ and/or $N_X$ (and also $a_Y$ and $a_Z$), the number of times experiments $Y$ or $X$, respectively, can be done. Specifically, note that as $a_Z \to \infty$ and $N_Y \to \infty$, we have $\mathsf{E}_\infty [\tau_{\textup{\fontsize{6 pt}{10 pt}\selectfont 2E}}(a_Z D_{\sigma^Z}, 0, a_Y, N_X, N_Y)] \to \infty$ and, thus, $\textup{POR}_Z \to 0$.

\medskip

The expression for $\textup{POR}_Y$ also follows from the Renewal Reward Theorem because in the renewal cycles starting with statistic $D_n=0$ to the point where $D_n$ is again zero after being negative, the time spent using experiment $Y$ is equal to the time for a truncated (at $N_Y$) CUSUM statistic for experiment $Y$ to reach $0$ after starting and getting reflected at the negative level $a_Z D_{\sigma^Z}$. This time is statistically equal to $\min \{ \tau_{\textup{\fontsize{6 pt}{10 pt}\selectfont C}}^Y (\vert a_Z D_{\sigma^Z} \vert), N_Y \}$. Thus,
\begin{align*}
 \textup{POR}_Y (\tau_{\textup{\fontsize{6 pt}{10 pt}\selectfont 3E}}) = \frac{\mathsf{E}_\infty [\min \{ \tau_{\textup{\fontsize{6 pt}{10 pt}\selectfont C}}^Y (\vert a_Z D_{\sigma^Z} \vert), N_Y \}]}{\mathsf{E}_\infty [\sigma^Z] + \mathsf{E}_\infty [\tau_{\textup{\fontsize{6 pt}{10 pt}\selectfont 2E}}(a_Z D_{\sigma^Z}, 0, a_Y, N_X, N_Y)]}.
\end{align*}
To control $\textup{POR}_Y$, first note that as $N_X \to 0$,
\begin{align*}
 \textup{POR}_Y (\tau_{\textup{\fontsize{6 pt}{10 pt}\selectfont 3E}}) &= \frac{\mathsf{E}_\infty [\min \{ \tau_{\textup{\fontsize{6 pt}{10 pt}\selectfont C}}^Y (\vert a_Z D_{\sigma^Z} \vert), N_Y \}]}{\mathsf{E}_\infty [\sigma^Z] + \mathsf{E}_\infty [\tau_{\textup{\fontsize{6 pt}{10 pt}\selectfont 2E}}(a_Z D_{\sigma^Z}, 0, a_Y, N_X, N_Y)]} \\
& \to \frac{\mathsf{E}_\infty [\min \{ \tau_{\textup{\fontsize{6 pt}{10 pt}\selectfont C}}^Y (\vert a_Z D_{\sigma^Z} \vert), N_Y \}]}{\mathsf{E}_\infty [\sigma^Z] + \mathsf{E}_\infty [\min \{ \tau_{\textup{\fontsize{6 pt}{10 pt}\selectfont C}}^Y (\vert a_Z D_{\sigma^Z} \vert), N_Y \}]}.
\end{align*}
Thus, one way to increase $\textup{POR}_Y$ is to bring $N_X$ to $0$ and increase $a_Z$ and $N_Y$. On the other hand, to decrease $\textup{POR}_Y (\tau_{\textup{\fontsize{6 pt}{10 pt}\selectfont 3E}})$, we increase $N_X$ and $a_Y$. 

\medskip

Finally, we have
$$ \textup{POR}_X (\tau_{\textup{\fontsize{6 pt}{10 pt}\selectfont 3E}}) = 1 - \textup{POR}_Y (\tau_{\textup{\fontsize{6 pt}{10 pt}\selectfont 2E}}) - \textup{POR}_Z (\tau_{\textup{\fontsize{6 pt}{10 pt}\selectfont 2E}}). $$
Note that the expected number of observations of experiment $X$ in a cycle can be written as
\begin{align*}
 \mathsf{E}_\infty [\tau_{\textup{\fontsize{6 pt}{10 pt}\selectfont 2E}}(a_Z D_{\sigma^Z}, 0, a_Y, N_X, N_Y)]  - \mathsf{E}_\infty [\min \{ \tau_{\textup{\fontsize{6 pt}{10 pt}\selectfont C}}^Y (\vert a_Z D_{\sigma^Z} \vert), N_Y \}].
\end{align*}

\item For the delay analysis of the 3E-CUSUM algorithm, suppose the change occurs at time $\nu = 2$. Then the worst possible realization from experiment $Z$ at time $1$ would create a potentially infinite amount of undershoot $D_{\sigma^Z}$. Thus, the only way for the statistic $D_n$ to come back to $0$ would be to exhaust all the $N_Y$ samples from the experiment $Y$. Starting time 2, each time an observation from experiment $Y$ is chosen, there is a potential $N_X$ number of observations one can collect from experiment $X$. Thus, one can bound the number of samples before $D_n$ comes above zero by 
$$ N_Y + N_Y N_X. $$
Once the statistic $D_n$ comes above $0$, it is reset to $0$, and the delay from that point onward is 
$$ \mathsf{E}_1 [\tau_{\textup{\fontsize{6 pt}{10 pt}\selectfont 3E}}]. $$
A similar argument for other change points gives us 
$$ \textup{WADD} (\tau_{\textup{\fontsize{6 pt}{10 pt}\selectfont 3E}}) \leq \mathsf{E}_1 [\tau_{\textup{\fontsize{6 pt}{10 pt}\selectfont 3E}}] + N_Y + N_Y N_X. $$
Using an argument similar to that used in the case of the 2E-CUSUM algorithm, we have
\begin{align*}
\textup{WADD} (\tau_{\textup{\fontsize{6 pt}{10 pt}\selectfont 3E}})
& \leq \mathsf{E}_1 [\tau_{\textup{\fontsize{6 pt}{10 pt}\selectfont 3E}}] + N_Y + N_Y N_X \\
& = (\mathsf{E}_1 [\tau_{\textup{\fontsize{6 pt}{10 pt}\selectfont C}}^Z] + K) + N_Y + N_Y N_X \\
& = \textup{WADD} (\tau_{\textup{\fontsize{6 pt}{10 pt}\selectfont C}}^Z) + K + N_Y + N_Y N_X.
\end{align*}
Here $K$ is a constant that does not depend on the threshold $A$. 
\end{enumerate}
\end{proof}

\end{document}
